\scanpage{118}
% archive sheet number: 117
En mettant ensemble les équations pertinentes : --
\margin{$z_i^2=2\unsure{\zeta_i|\xi_i}z_i-(a_2+2\zeta_j\zeta_k)\ \Longleftrightarrow$}
\[
 \left\{\begin{aligned}
 &(z_i-\zeta_i)^2=-\tfrac1{\zeta_i}(a_1\zeta_i^2+\unsure{2|L}a_2\zeta_i+3a_3)\quad [=-(a_1\zeta_i+2a_2+3\zeta_j\zeta_k)\\
 &{z'_i}^{\,2}=z_i^3+a_1z_i^2+a_2z_i+a_3\quad [=(z_i-\zeta_i)(z_i-\zeta_j)(z_i-\zeta_k)]\\
 &\det\begin{pmatrix}1&z_1&z'_1\\1&z_2&z'_2\\1&z_3&z'_3\end{pmatrix}=0\quad\text{i.e.}\quad
 \underbrace{(z_2z'_3-z_3z'_2)}_{\psi_1}+\underbrace{(z_3z'_1-z_1z'_3)}_{\psi_2}+\underbrace{(z_1z'_2-z_2z'_1)}_{\psi_3}=0
 \end{aligned}\right. \tag*{(71)}
\]
\margin{\uline{NB.} $[\psi_i^\omega=z_jz'_k-z_kz'_j\quad j=\omega i,\ k=\omega^2i=\omega j$}
qui expriment conjointement \deleted{\unsure{les|la}} donnée de $(x_i)_{i\in\unsure{I|J}}\in\unsure{\mathbb{E}|E}(k)^I$ satisfaisant (67) (69),
\[
 2x_i=\varepsilon_i\ (\in\phantom{E}_2E(k)^*),\qquad \sum x_i=0 \tag*{(72)}
\]
via
\[
 z_i=\wp(x_i),\quad z'_i=\wp'(x_i). \tag*{(73)}
\]

Je voudrais, à une telle donnée$_2$ $(x_i)_{i\in I}$, exprimée calculatoirement (via le choix d'un $\delta\in t_E^*$ et $\wp\in\unsure{\mathcal F_2|\mathscr F_2}$ associés) par (71), associer \inserted{canoniquement} un \uline{épinglage de Legendre de $E$}, de façon fonctorielle en $E$ et compatible avec les changem$^{\mathrm{ts}}$ de corps de base (entre autres\ldots). Je vais pour ceci supposer d'abord fixé un
\[
 \unsure{\mathbb{i}|\boldsymbol{i}|i|\dot{\imath}}\in\mu_4^*(k), \tag*{(74)}
\]
de sorte que, via ce choix, la première donnée $\alpha$ d'un épinglage de Legendre de $E$ s'interprète comme un élément
\[
 \omega\in\underline{\omega}=\underline{\omega}_I\qquad
 \begin{aligned}&\text{un ordre circulaire sur}\\ &I=\phantom{E}_2E(k)^*=\{\varepsilon_i\},\end{aligned}
 \tag*{(75)}
\]
% slanted
\margin{$\alpha=\mathbb{i}\wedge\omega$}
de sorte que la deuxième donnée $\kappa$ s'interprète comme un isomorphisme :
\[
 \kappa:\underline{t}_E\xrightarrow{\sim}\mathcal O_{\Sigma_I,t}(1)^{\unsure{\sharp|\natural}}
 \xrightarrow[\ ]{\inserted{\text{grâce au choix de }\mathbb{i}}\,\sim}\mathcal O_t(1),\tag*{(75)}
\]
